% GENERATED FILE. Edit canonical lessons, then run npm run generate:books.
% canonical-source-hash: fnv1a64:d9c130f5e50a68b7
% canonical-lessons: TA-C07-numbers-1-5, TA-C07-numbers-1-5-family, TA-C07-numbers-6-10, TA-W03-pulli-vanakkam, TA-C07-numbers-6-10-family

\chapter{Numbers One to Ten}
\label{ch:ta-numbers-one-to-ten}
\hlchaptermodality{7}

\begin{chapteropening}
\textbf{By the end of this chapter:} \emph{I can count one to ten in Tamil and say which of those numbers are inherited across the Dravidian family rather than borrowed.}

Pair Tamil pattu with Kannada hattu, apply the same p$\to$h softening to pāl/hālu, and identify ēḻu as one inherited seven.
\end{chapteropening}

\section[\ta{ஒன்று} \ta{இரண்டு} \ta{மூன்று} \ta{நான்கு} \ta{ஐந்து}]{\ta{ஒன்று}, \ta{இரண்டு}, \ta{மூன்று}, \ta{நான்கு}, \ta{ஐந்து} — one to five}

\label{lesson:TA-C07-numbers-1-5}



\begin{quote}
Five words — and the same point \emph{vaṇakkam} made, told again in numbers: Tamil counts with \textbf{its own} words, not borrowed ones.
\end{quote}

\subsection*{You'll want to know: The five}
\noindent
\begin{tabularx}{\linewidth}{@{}>{\raggedright\arraybackslash}X>{\raggedright\arraybackslash}X>{\raggedright\arraybackslash}X>{\raggedright\arraybackslash}X@{}}
\toprule
\textbf{} & \textbf{numeral} & \textbf{word} & \textbf{said} \\
\midrule
1 & \textbf{\ta{௧}} & \textbf{\ta{ஒன்று}} & \emph{oṉṟu} \\
2 & \textbf{\ta{௨}} & \textbf{\ta{இரண்டு}} & \emph{iraṇṭu} \\
3 & \textbf{\ta{௩}} & \textbf{\ta{மூன்று}} & \emph{mūṉṟu} \\
4 & \textbf{\ta{௪}} & \textbf{\ta{நான்கு}} & \emph{nāṉku} \\
5 & \textbf{\ta{௫}} & \textbf{\ta{ஐந்து}} & \emph{aintu} \\
\bottomrule
\end{tabularx}

Two writing systems sit side by side here: the \textbf{word} (spelled out in letters) and the \textbf{numeral} (a single symbol — \ta{௧} \ta{௨} \ta{௩} \ta{௪} \ta{௫}, Tamil's own digits, as distinct from 1 2 3 4 5 as they are from Roman I II III). You'll meet the digits again whenever a price or a page number is written in Tamil.

If you're doing the writing track, you can already read pieces of these: \textbf{\ta{ந}} and \textbf{\ta{ன}} (the \emph{n}-family), \textbf{\ta{ற}} the hard \emph{r} inside \emph{mūṉṟu}, and the \textbf{puḷḷi} dot that stacks \emph{\ta{ன்ற}} into the single cluster you hear in \textbf{\ta{ஒன்று}} and \textbf{\ta{மூன்று}}.

\subsection*{Your turn}
Say these aloud:

\begin{itemize}
\raggedright
  \item \textquotedblleft{}oṉṟu, iraṇṭu, mūṉṟu, nāṉku, aintu\textquotedblright{}
  \item read the digits aloud — \ta{௧} \ta{௨} \ta{௩} \ta{௪} \ta{௫}
\end{itemize}

\begin{tcolorbox}[breakable,colback=teal!4,colframe=teal!35!black,title={Before you move on}]
Count to five in Tamil. (\emph{Oṉṟu, iraṇṭu, mūṉṟu, nāṉku, aintu}.) Are Tamil's numbers borrowed from Sanskrit, like Hindi's? (\textbf{No} — they're native Dravidian.) Which two writing systems sit side by side? (Spelled-out \textbf{words} and Tamil's own \textbf{numerals}.) Next: compare these native forms across the family.
\end{tcolorbox}
\section[iraṇṭu mūṉṟu nāṉku aintu]{Two to five — one family in four regional coats}

\label{lesson:TA-C07-numbers-1-5-family}



\begin{quote}
Hindi's number forms are Sanskrit worn smooth. Tamil's are native Dravidian, so their sister-language shapes reveal the family directly.
\end{quote}

\begin{cousinweb}
\noindent
\begin{tabularx}{\linewidth}{@{}>{\raggedright\arraybackslash}X>{\raggedright\arraybackslash}X>{\raggedright\arraybackslash}X>{\raggedright\arraybackslash}X>{\raggedright\arraybackslash}X@{}}
\toprule
\textbf{} & \textbf{Tamil} & \textbf{Kannada} & \textbf{Telugu} & \textbf{Malayalam} \\
\midrule
2 & \emph{iraṇṭu} & \emph{eraḍu} & \emph{reṇḍu} & \emph{raṇṭu} \\
3 & \emph{mūṉṟu} & \emph{mūru} & \emph{mūḍu} & \emph{mūnnu} \\
4 & \emph{nāṉku} & \emph{nālku} & \emph{nālugu} & \emph{nālu} \\
5 & \emph{aintu} & \emph{aidu} & \emph{aidu} & \emph{añcu} \\
\bottomrule
\end{tabularx}

Each row is one word in four regional coats. That shared skeleton is why one Dravidian language gives you a running start on another.

\textbf{One} is less tidy: Tamil \emph{oṉṟu}, Kannada \emph{ondu}, and Malayalam \emph{onnu} share old \emph{oṉ-}, while Telugu uses the separate formation \emph{okaṭi}. The family agrees more neatly on 2–5 than on 1.
\end{cousinweb}

\begin{sounds}
In \emph{oṉṟu} and \emph{mūṉṟu}, \textbf{\ta{ன்ற}} is alveolar \emph{n} running into hard \emph{ṟ}, almost one knock: \emph{on-ru, muun-ru}. Puḷḷi on \textbf{\ta{ன்}} removes its vowel before the cluster.
\end{sounds}

\subsection*{Your turn}
\begin{itemize}
\raggedright
  \item \emph{Say it:} \textquotedblleft{}iraṇṭu · eraḍu · reṇḍu · raṇṭu\textquotedblright{}
  \item \emph{Say it:} \emph{oṉṟu} and \emph{mūṉṟu}, holding \textbf{\ta{ன்ற}} together
  \item \emph{Identify:} Telugu \emph{okaṭi} as the different \textquotedblleft{}one\textquotedblright{}
\end{itemize}

\begin{tcolorbox}[breakable,colback=teal!4,colframe=teal!35!black,title={Before you move on}]
Are these forms borrowed? (\textbf{No}, native Dravidian.) Which numbers stay tidiest across the family? (\textbf{Two through five}.) Which language forms \textquotedblleft{}one\textquotedblright{} differently? (\textbf{Telugu}, \emph{okaṭi}.)
\end{tcolorbox}
\section[\ta{ஆறு} \ta{ஏழு} \ta{எட்டு} \ta{ஒன்பது} \ta{பத்து}]{\ta{ஆறு}, \ta{ஏழு}, \ta{எட்டு}, \ta{ஒன்பது}, \ta{பத்து} — six to ten}

\label{lesson:TA-C07-numbers-6-10}



\begin{quote}
Five more — and one of them is written with the single letter Tamil is most proud of.
\end{quote}

\subsection*{You'll want to know: The five}
\noindent
\begin{tabularx}{\linewidth}{@{}>{\raggedright\arraybackslash}X>{\raggedright\arraybackslash}X>{\raggedright\arraybackslash}X>{\raggedright\arraybackslash}X@{}}
\toprule
\textbf{} & \textbf{numeral} & \textbf{word} & \textbf{said} \\
\midrule
6 & \textbf{\ta{௬}} & \textbf{\ta{ஆறு}} & \emph{āṟu} \\
7 & \textbf{\ta{௭}} & \textbf{\ta{ஏழு}} & \emph{ēḻu} \\
8 & \textbf{\ta{௮}} & \textbf{\ta{எட்டு}} & \emph{eṭṭu} \\
9 & \textbf{\ta{௯}} & \textbf{\ta{ஒன்பது}} & \emph{oṉpatu} \\
10 & \textbf{\ta{௰}} & \textbf{\ta{பத்து}} & \emph{pattu} \\
\bottomrule
\end{tabularx}

\subsection*{You'll want to know: Seven carries the letter that names the language}
Look at \textbf{\ta{ஏழு}} (\emph{ēḻu}, seven). Its middle letter is \textbf{\ta{ழ}} — the retroflex approximant, the sound English can only clumsily spell \emph{zh}. It is the rarest and most characteristic sound in Tamil, and it sits inside the language's \textbf{own name}: \emph{Tamiḻ} (\ta{தமிழ்}) ends in this very letter. So \emph{seven} is a free daily rehearsal of the sound that says who the language is. Curl the tongue back, don't touch, and voice it: \emph{ē-ḻu}.

(You met its cousin \emph{\ta{ற}} — the hard \emph{ṟ} — back in \emph{\ta{ஆறு}} / \emph{āṟu}, six. Tamil keeps these apart carefully: \textbf{\ta{ஆறு}} \emph{āṟu} and \textbf{\ta{ஏழு}} \emph{ēḻu} end in two different sounds English would both call \textquotedblleft{}r/zh\textquotedblright{}.)

\begin{grammarlens}[title={Grammar Lens: Nine is built from words you already have}]
\textbf{\ta{ஒன்பது}} (\emph{oṉpatu}, nine) is not a fresh word to memorise — it is two you already know, joined:

\begin{quote}
\textbf{\ta{ஒன்}} (\emph{oṉ-}, \textquotedblleft{}one,\textquotedblright{} from last lesson) + \textbf{\ta{பது}} (\emph{patu}, \textquotedblleft{}ten\textquotedblright{}) $\to$ \textbf{\ta{ஒன்பது}} = \textquotedblleft{}one‑to‑ten\textquotedblright{} — one short of ten.
\end{quote}

And ten itself is \textbf{\ta{பத்து}} (\emph{pattu}). Say \emph{oṉpatu} then \emph{pattu} and you can hear nine leaning on ten. That is a common Dravidian habit — nine described as one-away-from-ten rather than named outright.
\end{grammarlens}

\subsection*{Your turn}
Say these aloud:

\begin{itemize}
\raggedright
  \item \textquotedblleft{}āṟu, ēḻu, eṭṭu, oṉpatu, pattu\textquotedblright{}
  \item seven, feeling the \ta{ழ} — \textquotedblleft{}ē‑ḻu\textquotedblright{}, the sound in \emph{Tamiḻ}
  \item nine taken apart — \textquotedblleft{}oṉ + patu $\to$ oṉpatu\textquotedblright{}, one short of ten
  \item read the digits — \ta{௬} \ta{௭} \ta{௮} \ta{௯} \ta{௰}
\end{itemize}

\begin{tcolorbox}[breakable,colback=teal!4,colframe=teal!35!black,title={Before you move on}]
Count six to ten in Tamil. (\emph{Āṟu, ēḻu, eṭṭu, oṉpatu, pattu}.) Which number is written with \textbf{\ta{ழ}}, and where else does that letter appear? (\textbf{Seven}, \emph{ēḻu} — the same letter that ends \emph{Tamiḻ}.) How is \textbf{nine} built? (\emph{Oṉ} \textquotedblleft{}one\textquotedblright{} + \emph{patu} \textquotedblleft{}ten\textquotedblright{} $\to$ \emph{oṉpatu}, one short of ten.) Next: the puḷḷi — the dot that takes a consonant's built-in vowel away.
\end{tcolorbox}
\section[puḷḷi]{\ta{்} — the dot that leaves both letters visible}

\label{lesson:TA-W03-pulli-vanakkam}



\begin{quote}
Every consonant carries a built-in \textbf{a}. So how do you write a consonant with \textbf{no} vowel — the \emph{k} in \emph{vaṇakkam}, or the \emph{m} at the end?

One dot.
\end{quote}

\subsection*{Script you'll notice: The puḷḷi}
\begin{quote}
\textbf{\ta{க}} = \emph{ka}  $\to$  \textbf{\ta{க்}} = \emph{k}
\end{quote}

\textbf{\ta{புள்ளி}} (\emph{puḷḷi}) means simply \textbf{\textquotedblleft{}the dot.\textquotedblright{}} Written above a letter, it removes the inherent vowel. That's the whole rule.

\textbf{A standing note for this whole track.} These lessons quote real Tamil words and examples, and most of them contain letters the track has not taught you to draw yet — \emph{puḷḷi}'s own name uses \textbf{\ta{ப}}, \textbf{\ta{ள}} and the \textbf{\ta{ு}} sign, and the comparisons below reach for others again. \textbf{Read anything printed here; only draw the letters a lesson has actually broken into strokes.} The data file behind this track covers 11 letters so far and says so (\texttt{complete: false}) — it grows as the chapters need it.

\subsection*{Script you'll notice: What Tamil does \emph{not} do}
If you have done the Devanagari track, this is the moment to notice a real divergence.

In Devanagari, killing the vowel makes the bare consonant \textbf{fuse} with the next one into a squeezed \textbf{conjunct} — \dv{स्} + \dv{त} becomes \dv{स्त}, and the first letter loses its spine. The virama usually disappears into the ligature.

\textbf{Tamil doesn't do that.} The dot stays visible, and both letters keep their full shape:

\noindent
\begin{tabularx}{\linewidth}{@{}>{\raggedright\arraybackslash}X>{\raggedright\arraybackslash}X@{}}
\toprule
\textbf{} & \textbf{} \\
\midrule
Devanagari & \dv{स्} + \dv{त} $\to$ \textbf{\dv{स्त}} — a new fused shape \\
\textbf{Tamil} & \ta{க்} + \ta{க} $\to$ \textbf{\ta{க்க}} — two letters, dot intact \\
\bottomrule
\end{tabularx}

The consequence is arithmetic. Tamil's whole set is \textbf{247 characters}: 12 vowels + 18 consonants + the 216 consonant-vowel combinations + the \textbf{āytam \ta{ஃ}}. That is genuinely all of them. Devanagari's conjuncts run to many hundreds of ligature shapes on top of its letters, each to be learned or at least recognised.

(One honest exception: Tamil does have a few \textbf{ligatures borrowed with Sanskrit words} — \textbf{\ta{க்ஷ}} \emph{kṣa} and \textbf{\ta{ஸ்ரீ}} \emph{śrī}. They sit outside the 18 native consonants and outside the 247.)

\subsection*{Your turn}
Write these out:

\begin{itemize}
\raggedright
  \item \ta{க்} — \textquotedblleft{}\ta{க}, then the dot above\textquotedblright{}
  \item \ta{க்க} — two full letters, dot intact, \textbf{no fusing}
\end{itemize}

\begin{tcolorbox}[breakable,colback=teal!4,colframe=teal!35!black,title={Before you move on}]
What does the \textbf{puḷḷi} do, and what does the word mean? (Removes the \textbf{inherent vowel}; \emph{puḷḷi} is simply \textquotedblleft{}\textbf{the dot}\textquotedblright{}.) What does Devanagari do that Tamil doesn't? (\textbf{Fuses} the letters into a \textbf{conjunct} — Tamil keeps both shapes and the dot.) What does Tamil get in exchange? (A \textbf{much smaller character set} — 247, with \textbf{almost} no ligatures: only the borrowed \ta{க்ஷ} and \ta{ஸ்ரீ}.) Next: six, seven and ten across the family.
\end{tcolorbox}
\section[āṟu ēḻu pattu]{Six to ten — family resemblances and one Kannada shift}

\label{lesson:TA-C07-numbers-6-10-family}



\begin{quote}
The Tamil forms are learned. Now use their sister shapes as evidence for inheritance and regular sound change.
\end{quote}

\begin{cousinweb}
Six and seven remain close: Tamil \emph{āṟu}, Kannada \emph{āru}, Malayalam \emph{āṟu}; Tamil \emph{ēḻu}, Kannada \emph{ēḷu}, Malayalam \emph{ēḻu}. Tamil and Malayalam also keep eight and nine nearly identical: \emph{eṭṭu} and \emph{oṉpatu/onpatu}.

Telugu wanders on eight and nine with \emph{enimidi} and \emph{tommidi}, so those forms are less reliable as family-wide recognition anchors.
\end{cousinweb}

\begin{cousinweb}
Tamil and Malayalam say \textbf{pattu}, while Kannada says \textbf{hattu}. The initial \emph{p $\to$ h} is a regular Kannada shift, not a random replacement. Compare Tamil \emph{pāl}, \textquotedblleft{}milk,\textquotedblright{} with Kannada \emph{hālu}.

A correspondence lets differences prove relationship: repeated \emph{p/h} pairs are stronger evidence than one accidental resemblance.
\end{cousinweb}

\subsection*{Your turn}
\begin{itemize}
\raggedright
  \item \emph{Say it:} Tamil \emph{pattu} · Kannada \emph{hattu}
  \item \emph{Say it:} Tamil \emph{pāl} · Kannada \emph{hālu}
  \item \emph{Identify:} \emph{ēḻu / ēḷu / ēḻu} as one inherited seven
\end{itemize}

\begin{tcolorbox}[breakable,colback=teal!4,colframe=teal!35!black,title={Before you move on}]
What happens to Tamil \emph{pattu} in Kannada? (\emph{P} regularly softens to \emph{\textbf{h}}: \emph{hattu}.) Which second pair confirms the pattern? (\emph{Pāl / hālu}, \textquotedblleft{}milk.\textquotedblright{}) Which language diverges most on eight and nine? (\textbf{Telugu}.)
\end{tcolorbox}
